\documentclass[11pt,a4paper]{article}

\usepackage[utf8]{inputenc}
\usepackage[T1]{fontenc}
\usepackage{lmodern}
\usepackage{amsmath,amssymb,amsthm,mathtools}
\usepackage{geometry}
\geometry{margin=2.7cm}
\usepackage{hyperref}
\usepackage{enumitem}
\usepackage{xcolor}
\usepackage{booktabs}
\usepackage{tcolorbox}
\tcbuselibrary{breakable,skins}
\usepackage{microtype}

\newtheorem{theorem}{Theorem}[section]
\newtheorem{lemma}[theorem]{Lemma}
\newtheorem{proposition}[theorem]{Proposition}
\newtheorem{corollary}[theorem]{Corollary}
\newtheorem{conjecture}[theorem]{Conjecture}
\theoremstyle{definition}
\newtheorem{definition}[theorem]{Definition}
\newtheorem{remark}[theorem]{Remark}

\newcommand{\Q}{\mathbb{Q}}
\newcommand{\Z}{\mathbb{Z}}
\newcommand{\R}{\mathbb{R}}
\newcommand{\C}{\mathbb{C}}
\newcommand{\F}{\mathbb{F}}
\newcommand{\A}{\mathbb{A}}
\newcommand{\OK}{\mathcal{O}_K}
\newcommand{\Cphi}{\mathcal{C}_{\!\varphi}}
\newcommand{\Vmin}{V_{\min}}
\newcommand{\GL}{\mathrm{GL}}

% status tags
\newcommand{\TAGthm}{{\small\textsf{[\textcolor{green!45!black}{theorem}]}}}
\newcommand{\TAGver}{{\small\textsf{[\textcolor{blue!60!black}{verified}]}}}
\newcommand{\TAGev}{{\small\textsf{[\textcolor{orange!80!black}{evidence}]}}}
\newcommand{\TAGconj}{{\small\textsf{[\textcolor{purple}{conjecture}]}}}
\newcommand{\TAGopen}{{\small\textsf{[\textcolor{red!70!black}{open}]}}}

\hypersetup{colorlinks=true, linkcolor=blue!55!black,
  citecolor=blue!55!black, urlcolor=blue!55!black}

\title{\textbf{The $V_{600}$ Substrate at the Riemann--Hypothesis Frontier}\\
\large A verified local realisation, a corrected reduction,\\ and an honest map of the global gap}
\author{Lee Smart \\ \small Institute of Vibrational Field Dynamics}
\date{2026-05-30 \quad v1.0}

\begin{document}
\maketitle

\begin{center}
\fbox{\parbox{0.93\textwidth}{\centering\small
\textsc{Pre-peer-review open research preprint.} \textbf{No proof of the
Riemann Hypothesis is claimed.} Every claim below is tagged
\textsf{[theorem]}, \textsf{[verified]} (computed, with a stated correctness
gate), \textsf{[evidence]} (numerics consistent with, not proving),
\textsf{[conjecture]}, or \textsf{[open]}. All computations are reproducible
from the accompanying \texttt{route\_b/} and \texttt{sims/} code.}}
\end{center}

\begin{abstract}
We report the honest standing of the $V_{600}$ closure substrate at the
Riemann--Hypothesis (RH) frontier. \textbf{Main verified result:} the
icosian ring --- the maximal order of the definite quaternion algebra
$(-1,-1/\Q(\sqrt5))$ whose unit group is the $600$-cell --- computes, by its
own intrinsic Brandt/Hecke action at level $\mathfrak{p}_{31}$ (norm $31$),
the genuine cuspidal Hilbert modular newform over $\Q(\sqrt5)$ of that level:
its Hecke eigenvalues match the independently point-counted newform exactly
on all $11$ prime ideals of norm $\le 41$, including signs and split-prime
pairings, out-of-sample, with no fitting. This is a non-circular
substrate$\leftrightarrow$arithmetic identification, replacing a hardcoded
map a prior audit exposed. \textbf{Corrected scope:} we withdraw an earlier
claim that admissibility plus the explicit formula yields RH; they certify a
genuine automorphic L-function, whose RH is equivalent to GRH and remains
open. \textbf{Geometric map:} the substrate realises the \emph{local}
critical line (the per-prime Satake circle $|\alpha_\mathfrak{q}|=\sqrt{N\mathfrak{q}}$,
forced by temperedness) but not the \emph{global} one; we explain why
local$\,\to\,$global is the genuine wall and name the missing object (a
Hilbert--Pólya operator on the adelic arena). \textbf{Frontier computation:}
we compute the Connes/Weil positivity functional for this L-function from the
substrate's own Satake data, gated against the Riemann $\zeta$ zeros to
$0.01\%$; it is positive on every validated test function --- evidence
consistent with RH, not a proof. \textbf{Locating the operator:} the zeros
are GUE (verified: level repulsion; a quasicrystal dual to the primes), and
the Dyson/division-algebra structure pins the missing Hilbert--Pólya operator
to the $\C$ (unitary, $\beta=2$, time-reversal-broken) class --- one rung
from the substrate's $\mathbb{H}$ ($\beta=4$). The $\mathbb{H}\to\C$
transition is complexification ($\mathbb{H}\otimes\C=M_2(\C)$,
$2I\subset\mathrm{SU}(2)$, verified) realised by Jacquet--Langlands; the
finite-place part of the operator (the commutative Hecke scaling algebra) is
built and verified $\C$-class, leaving the archimedean self-adjoint
realisation --- equivalently Weil positivity for all test functions, i.e.\
RH --- as the single open step.
\end{abstract}

\tableofcontents

\section{Stance}
This paper does not attempt to prove RH. It records, with reproducible
computations and explicit status tags, exactly how far a specific geometric
substrate reaches into the problem, where it stops, and why. The discipline
throughout: the mathematics stays fully mathematical; no interpretive,
observer-, or consciousness-based reading is load-bearing or implied.

\section{The substrate and the local critical line}\label{sec:local}

\subsection{The icosian ring}
Let $K=\Q(\sqrt5)$, $\OK=\Z[\varphi]$, $\varphi=(1+\sqrt5)/2$. Let
$B=(-1,-1/K)$ be the Hamilton quaternions over $K$, ramified at both real
places (definite). Its maximal order is the \emph{icosian ring} $\mathcal{I}$,
whose group of units is the binary icosahedral group $2I$ of order $120$ ---
the $120$ vertices of the $600$-cell $V_{600}$.

\begin{proposition}[\TAGver]\label{prop:icosian}
The $120$ unit icosians all have reduced norm $1$ and are closed under
multiplication ($14400/14400$ products verified in
\texttt{route\_b/icosian.py}), forming $2I$.
\end{proposition}

\subsection{The two critical lines}
A central clarification of this work is that ``the critical line'' refers to
two different objects:
\begin{enumerate}[label=(\arabic*)]
\item \textbf{Local:} for each prime $\mathfrak{q}$, the zeros of the local
factor $L_\mathfrak{q}(s)$. The Satake parameters satisfy
$\alpha_\mathfrak{q}=\sqrt{N\mathfrak{q}}\,e^{i\theta_\mathfrak{q}}$ on the
circle $|\alpha_\mathfrak{q}|=\sqrt{N\mathfrak{q}}$ --- temperedness /
Ramanujan --- and consequently \emph{each} local factor's zeros lie exactly
on $\Re(s)=1/2$.
\item \textbf{Global:} the zeros of the full product $L(s)=\prod_\mathfrak{q}
L_\mathfrak{q}(s)$. That these lie on $\Re(s)=1/2$ is RH.
\end{enumerate}

\begin{proposition}[\TAGver, \texttt{route\_b/satake\_circle.py}]
For the substrate-computed eigenvalues (\S\ref{sec:circle}), every Satake
parameter lies on its circle and every local zero lies on $\Re(s)=1/2$;
e.g.\ at $N\mathfrak{q}=4,5,9,11$ the angles are
$138.6^\circ, 116.6^\circ, 70.5^\circ, 52.9^\circ/127.1^\circ$.
\end{proposition}

This is the number-field shadow of the function-field RH that Deligne proved
geometrically (Frobenius weights on $|\alpha|=\sqrt q$). The substrate
realises the local critical line; the global one is the open problem.

\subsection{Why the Hecke structure reaches the L-function}\label{sec:tilings}
The substrate carries three independent extension structures, and only one
reaches the L-function's zeros:
\begin{itemize}[leftmargin=1.4em]
\item \textbf{$W(H_4)$ geometric} (order $14{,}400$): extends a statement from
one $V_{600}$ vertex to all $120$ --- \emph{within} the substrate.
\item \textbf{Cascade scale} ($\mathrm{Pentagon}\!\to\!\mathrm{Icos}\!\to\!
V_{600}\!\to\!E_8\!\to\!\cdots$): propagates patterns up the hierarchy ---
\emph{within} the substrate.
\item \textbf{Hecke multiplicative}: the Hecke operators at fixed level are
generated by the prime operators, with $T_\mathfrak{m}T_\mathfrak{n}=
T_{\mathfrak{m}\mathfrak{n}}$ for coprime ideals and a prime-power recursion.
This is the only one of the three that reaches the global L-function.
\end{itemize}

\begin{theorem}[Multiplicativity tiling, \TAGthm]\label{thm:tiling}
Let $\{a_\mathfrak{p}\}_{N\mathfrak{p}\le P}$ be the eigenvalues of a commuting
family on a finite-dimensional space, satisfying Ramanujan
$|a_\mathfrak{p}|\le2\sqrt{N\mathfrak{p}}$, the prime-power recursion
$a_{\mathfrak{p}^{k+1}}=a_\mathfrak{p}a_{\mathfrak{p}^k}-N\mathfrak{p}\,
a_{\mathfrak{p}^{k-1}}$, and $a_{\mathfrak{a}\mathfrak{b}}=a_\mathfrak{a}
a_\mathfrak{b}$ for coprime $\mathfrak{a},\mathfrak{b}$. Then the full
coefficient system $\{a_\mathfrak{n}\}$ is determined by the finite prime data,
and
$$L(s,\psi)=\sum_\mathfrak{n}\frac{a_\mathfrak{n}}{N\mathfrak{n}^{s}}
=\prod_\mathfrak{p}\bigl(1-a_\mathfrak{p}N\mathfrak{p}^{-s}
+N\mathfrak{p}^{1-2s}\bigr)^{-1}$$
converges absolutely for $\Re s>3/2$ (Euler product unconditional).
\end{theorem}

Hence the circle test (\S\ref{sec:circle}) is decisive: matching the
substrate's \emph{finite} prime eigenvalues to the genuine newform pins the
\emph{entire} L-function by tiling. It does not, by itself, place the zeros ---
that is RH (\S\ref{sec:correct}).

\section{The circle test: the substrate computes a genuine newform}\label{sec:circle}

The target is the cuspidal Hilbert modular newform over $K$ at level
$\mathfrak{p}_{31}$ (norm $31$; the prime $31\equiv1\!\pmod5$ splits in $K$).
This is the level at which the first cuspidal newform over $\Q(\sqrt5)$
appears \cite{Dembele}, corresponding to the elliptic curve $31.1$-$a$.

\subsection{The genuine target $H$ (independent of the substrate)}
\begin{proposition}[\TAGver, \texttt{sims/sim\_genuine\_eigenvalues.py}]
\label{prop:target}
The Hecke eigenvalues $a_\mathfrak{q}=N\mathfrak{q}+1-\#E(\F_\mathfrak{q})$ of
the curve $31.1$-$a1$ ($y^2+xy+\varphi y=x^3+(\varphi+1)x^2+\varphi x$),
computed by point counting (no eigenvalue table used), for all $44$ good
prime ideals of norm $\le200$, satisfy:
the Ramanujan bound $|a_\mathfrak{q}|\le 2\sqrt{N\mathfrak{q}}$ (cuspidal),
and the torsion divisibility $8\mid \#E(\F_\mathfrak{q})$ (the curve has
$\Z/8\Z$ torsion --- an independent correctness proof needing no external
table).
\end{proposition}

LMFDB confirms a single isogeny class at norm $31$: the cuspidal space is
$1$-dimensional, not $26$, settling an earlier dimension question.

\subsection{The substrate side $S$ (from the icosians' own arithmetic)}
Because $\mathcal{I}$ has class number $1$, the level-$\mathfrak{p}_{31}$
modular space is $\C[\,2I\backslash \mathbb{P}^1(\F_{31})\,]$ (Voight
\cite{Voight}, Ch.\ 41; Dembl\'e \cite{Dembele}).

\begin{proposition}[\TAGver, \texttt{route\_b/step2\_eichler.py}]
Under the local splitting $B\otimes\F_{31}\cong M_2(\F_{31})$, the units $2I$
reduce to $A_5\subset\mathrm{PGL}_2(\F_{31})$ (order $60$), acting on
$\mathbb{P}^1(\F_{31})$ ($32$ points) with orbit sizes $\{12,20\}$. Hence
$\dim = 2 = 1$ Eisenstein $+\,1$ cuspidal, matching Proposition
\ref{prop:target}.
\end{proposition}

The Hecke operator $T_\mathfrak{q}$ is the sum of the norm-$\mathfrak{q}$
quaternions (enumerated by the verified rank-$8$ short-vector routine,
\texttt{route\_b/short\_vectors.py}, gated by $r(2)=120$) acting on the two
orbits; the cuspidal eigenvalue is $a_\mathfrak{q}=\mathrm{tr}(T_\mathfrak{q})-(N\mathfrak{q}+1)$.

\begin{theorem}[Circle closure, \TAGver, \texttt{route\_b/step3\_4\_hecke.py}]
\label{thm:circle}
The substrate cuspidal eigenvalues $S$ equal the genuine target $H$ on every
prime ideal of norm $\le41$, computed with \emph{no fitting} and \emph{no
modular input}:
\begin{center}
\begin{tabular}{lccccccc}
\toprule
$N\mathfrak{q}$ & $4$ & $5$ & $9$ & $11$ & $19$ & $29$ & $41$\\
\midrule
substrate $a_\mathfrak{q}$ & $-3$ & $-2$ & $2$ & $\{-4,4\}$ & $\{-4,4\}$ & $\{-2,-2\}$ & $\{-6,-6\}$\\
genuine $H$ & $-3$ & $-2$ & $2$ & $\{-4,4\}$ & $\{-4,4\}$ & $\{-2,-2\}$ & $\{-6,-6\}$\\
\bottomrule
\end{tabular}
\end{center}
The agreement is exact, includes signs and the split-prime pairings, and
holds \emph{out of sample} on $\{19,29,41\}$ (which influenced nothing). The
Eisenstein eigenvalue is $N\mathfrak{q}+1$ on the nose (row sums), and every
$a_\mathfrak{q}$ is Ramanujan-bounded.
\end{theorem}

\begin{remark}
This is the central verified result. It is, mathematically, \emph{expected}
--- $\mathcal{I}$ \emph{is} the maximal order of the relevant quaternion
algebra, so its Brandt action \emph{is} this newform by Jacquet--Langlands ---
which is exactly why it is the right, non-circular confirmation: the
substrate genuinely \emph{computes} (does not inject) the eigenvalues. It
replaces a hardcoded/fitted map that an internal provenance audit
(\texttt{CIRCLE\_TEST.md}) had exposed.
\end{remark}

\section{Corrected scope: admissibility does not give RH}\label{sec:correct}

The substrate eigenform passes the verifiable admissibility classes
\emph{genuinely}: multiplicativity $a_{\mathfrak{q}^2}=a_\mathfrak{q}^2-N\mathfrak{q}$
(checked: $a_{(2)^2}=5=(-3)^2-4$, $a_{(\sqrt5)^2}=-1=(-2)^2-5$,
\texttt{route\_b/admissibility.py}, \TAGver); Ramanujan (the Hasse bound,
\TAGthm); Sato--Tate (a theorem for this non-CM curve, \TAGthm); functional
equation and automorphy (\TAGthm). So $L(\psi,s)$ is a genuine degree-$4$
cuspidal automorphic L-function (conductor $5^2\cdot31=775$).

\begin{tcolorbox}[colback=red!2,colframe=red!50!black,arc=2mm,
  title={\textbf{Withdrawn overclaim}},fonttitle=\bfseries,breakable]
An earlier statement asserted that admissibility plus the Weil explicit
formula \emph{forces the zeros onto the critical line}. \textbf{This is
false and is withdrawn.} A genuine automorphic L-function with a valid
explicit formula does \emph{not} have its zeros constrained to the line ---
that is precisely RH for it, equivalent to GRH, and \TAGopen. The explicit
formula is an identity equivalent to information about the zeros; Ramanujan
bounds the prime side and Sato--Tate describes the eigenvalue distribution,
but \emph{neither pins the zeros}. (Re)established in
\texttt{route\_b/admissibility.py}.
\end{tcolorbox}

\section{Why local $\to$ global is the wall}\label{sec:wall}

The functional equation $\Lambda(s)=\Lambda(1-s)$ forces the zeros to be
\emph{symmetric about} $\Re(s)=1/2$ \TAGthm; it does \emph{not} force them
\emph{onto} it. A zero could be a symmetric pair $\tfrac12\pm\beta+i\gamma$
straddling the line; RH is the statement $\beta=0$ always. The structure
gives the \emph{axis}; RH is the \emph{on-ness}.

\begin{remark}[The global object is not a polytope]
The global L-function is the Euler product of all local factors, and its
zeros are \emph{emergent}: a product of factors each with zeros on a line has
its own new zeros, visible in no single factor. The object that glues all
places is the \emph{adele class space} $\A_\Q/\Q^\times$ --- whose compact
heart is the adelic solenoid, an \emph{inverse limit of circles}
(infinite-dimensional, noncommutative; Connes). It is not the $600$-cell,
not $S^3$, not any finite stack of cells. Climbing the cascade
($V_{600}\to E_8\to\cdots$) gives bigger \emph{local} objects, never the
global glue; the cascade limit \emph{is} the solenoid.
\end{remark}

\section{The adjoint that would close it}\label{sec:adjoint}

\begin{conjecture}[Hilbert--Pólya, \TAGconj]
There is a self-adjoint operator $\hat H$ on a Hilbert space with
$\mathrm{spec}(\hat H)=\{\gamma_n:\zeta(\tfrac12+i\gamma_n)=0\}$.
\end{conjecture}
Given such $\hat H$, self-adjointness gives real spectrum, hence RH \TAGthm.
The Weil explicit formula is already a trace formula in the exact shape of
the Selberg trace formula (primes $\leftrightarrow$ closed geodesics, $\log
p\leftrightarrow$ lengths, zeros $\leftrightarrow$ Laplace eigenvalues),
which is why $\hat H$ is believed to exist. Two written candidates: the
Berry--Keating Hamiltonian $xp$ (semiclassical count matches the zero count,
but $xp$ is not essentially self-adjoint), and Connes' operator on the adele
class space (a rigorous \emph{reformulation}, \S\ref{sec:connes}).

\begin{proposition}[The substrate's operator has the wrong spectrum, \TAGver]
$\Cphi=(12+\varphi^{-2})I-A_1$ on $V_{600}$ is self-adjoint and positive
definite, with spectrum
$\{0.382,2.67,5.91,9.38,12.38,14.38,14.85,15.38,16.09\}$ --- nine algebraic
numbers in $K$ (smallest $=\varphi^{-2}$, the witness $\Vmin$, multiplicity
$1$). It is the right \emph{type} (self-adjoint) but the wrong \emph{spectrum}
(finite, algebraic, local) for the infinite transcendental $\gamma_n$.
\end{proposition}

\section{Connes/Weil positivity, computed and gated}\label{sec:connes}

Connes' reformulation: RH for $L$ is equivalent to positivity of the Weil
distribution, $W(h)=(\text{archimedean})-(\text{prime})=\sum_\rho h(\gamma_\rho)\ge0$
for all $h$ of positive type. The prime term is built from the substrate's
verified Satake angles $\theta_\mathfrak{q}$; the archimedean term from the
gamma factor $\Gamma_\C(s)^2$ and conductor $775$.

\begin{proposition}[Validation gate, \TAGver, \texttt{route\_b/weil\_positivity.py}]
Run on the Riemann $\zeta$ function (known zeros), the identical machinery
reproduces $\sum_\rho h(\gamma_\rho)$ to relative error $1.0\%,\,0.01\%,\,0.00\%$
at Gaussian widths $\sigma=3,4,5$ --- confirming the constants, gamma factor,
and normalisation.
\end{proposition}

\begin{proposition}[Positivity holds for $L(\psi,s)$, \TAGev]
With the machinery gate-validated, the Weil/Connes functional for the
norm-$31$ newform, from the substrate's Satake data, is
$$ W(h) = +1.40,\ +3.70,\ +6.39 \quad\text{at } \sigma = 3,4,5, $$
positive on every validated test function.
\end{proposition}

\begin{tcolorbox}[colback=orange!3,colframe=orange!70!black,arc=2mm,
  title={\textbf{What this is and is not}},fonttitle=\bfseries,breakable]
This is \emph{evidence consistent with} RH for this L-function, computed from
substrate data and gated against $\zeta$. It is \textbf{not} a proof: Weil
positivity must hold for \emph{all} test functions of positive type, and we
test finitely many with the prime sum truncated ($N\mathfrak{q}\le150$). A
negative $W(h)$ would have \emph{disproved} RH; we found none. Proving
positivity for all $h$ \emph{is} RH, exactly as open as before.
\end{tcolorbox}

\section{The measurable bridge: quantum chaos}\label{sec:chaos}

The zeros connect to a \emph{measured} physical phenomenon. The unfolded
nearest-neighbour spacings of the Riemann zeros follow the \textbf{GUE}
random-matrix law --- the same law measured in heavy nuclei and chaotic
microwave cavities (Montgomery--Odlyzko; Rudnick--Sarnak).

\begin{proposition}[\TAGver, \texttt{route\_b/quantum\_chaos\_test.py}]
On $40$ zeros the spacings show level repulsion (no small gaps; fraction
$<0.3$ is $0.000$ vs Poisson $0.259$) and the histogram is $\sim$13$\times$
closer to GUE than to Poisson. The zeros statistically behave like a
quantum chaotic spectrum --- evidence for the spectral (Hilbert--Pólya)
picture, not a proof.
\end{proposition}

The same duality appears as a \emph{quasicrystal}: the diffraction
$D(t)=\sum_n\cos(\gamma_n t)$ has Bragg peaks exactly at $\log(p^k)$
(\TAGver, $11/11$ matched, \texttt{diffraction\_quasicrystal.py}) --- the
zeros (outside) diffract into the primes (inside), which is the explicit
formula seen as diffraction. Reconstructing the prime staircase from the
zeros layer-by-layer (\texttt{mirror\_layers.py}) shows each zero is a
frequency harmonic; off-line tuning grows the error as $x^{1/2+\beta}$, so
RH is the amplitude-balance condition.

\begin{remark}[Honest scope]
These confirm the structure (quantum-chaos class, quasicrystal duality);
peak \emph{locations} and spacing statistics hold regardless of RH. They are
texture evidence, not the operator.
\end{remark}

\section{The division-algebra structure}\label{sec:algebra}

The GUE spacing constant $\langle r\rangle\approx0.6027$ (Atas et al.) is a
function of the Dyson index $\beta$, and Dyson's threefold way identifies
$\beta$ with a division algebra:
$\beta{=}1{\leftrightarrow}\R$ (GOE), $\beta{=}2{\leftrightarrow}\C$ (GUE),
$\beta{=}4{\leftrightarrow}\mathbb{H}$ (GSE); physically these are the three
time-reversal classes, with $\beta{=}2$ (GUE) \textbf{time-reversal broken}.

\begin{proposition}[The class mismatch, \TAGver]
The substrate is $\mathbb{H}$ (quaternionic; natural class GSE, $\beta=4$):
its operator $A_1$ on $V_{600}$ is $93\%$ degenerate. The zeros are $\C$
(GUE, $\beta=2$; $\langle r\rangle=0.62$). Reaching the generator from the
substrate is a \emph{change of division algebra} $\mathbb{H}\to\C$, not a
deformation --- the sharpest statement of why the symmetric geometry cannot
itself be the generator.
\end{proposition}

\begin{proposition}[The $\mathbb{H}\to\C$ transition is complexification,
\TAGver, \texttt{route\_b/jacquet\_langlands\_trace.py}]
$\mathbb{H}\otimes\C\cong M_2(\C)$ (algebra iso; reduced norm $=$
determinant, to $2{\times}10^{-16}$), and the $120$ icosian units map into
$\mathrm{SU}(2)$ (to $2.5{\times}10^{-17}$): $2I\subset\mathrm{SU}(2)$, the
\emph{unitary} group. The Satake angles are $\mathrm{SU}(2)$ conjugacy
angles (Sato--Tate $=$ the $\mathrm{SU}(2)$ measure). Globally this is
Jacquet--Langlands ($B^\times\leftrightarrow\GL_2$, both becoming $M_2(\C)$);
the eigenvalue matching is the circle test (\S\ref{sec:circle}); the GUE high
zeros follow by Rudnick--Sarnak. \textbf{The class flip is complexification,
not geometry.}
\end{proposition}

\begin{proposition}[Symmetry type, \TAGver/\TAGthm,
\texttt{route\_b/katz\_sarnak.py}]
The $a_\mathfrak{q}$ are real $\Rightarrow$ $L$ is self-dual $\Rightarrow$
Katz--Sarnak \emph{family} symmetry is Orthogonal or Symplectic, not
Unitary; weight-$2$ $\GL_2$ (Iwaniec--Luo--Sarnak) $\Rightarrow$
\textbf{Orthogonal}, and rank $0$ $\Rightarrow$ $O^+$. So the full
$\R\C\mathbb{H}$ trichotomy appears in three roles:
substrate $=\mathbb{H}$, bulk zero statistics $=\C$ (GUE), family symmetry
$=\R$ (Orthogonal).
\end{proposition}

\section{Locating the operator: the adelic scaling}\label{sec:locate}

The conjectural operator is the generator of \emph{scale} (Berry--Keating
$xp$; Connes' idele action): RH $\Leftrightarrow$ scale is self-adjoint once
the primes discretise it (\texttt{SCALE\_AND\_SYMMETRY.md}). Decompose the
adelic scaling operator as \emph{finite places} (the Hecke algebra $=$ local
scalings) $\otimes$ \emph{archimedean} (the Mellin/dilation continuum).

\begin{proposition}[Finite part, \TAGver, \texttt{route\_b/adelic\_scaling.py}]
The Hecke operators $T_\mathfrak{q}$ on the Brandt module \emph{commute}
($\max\|[T_a,T_b]\|=0$) --- a commutative scaling algebra --- with joint
spectrum Eisenstein $(N\mathfrak{q}+1)\oplus$ cuspidal $(a_\mathfrak{q})$,
the latter unitary/Satake. So the adelic limit carries a $\C$-class
(unitary, $\beta=2$) scaling structure at the finite places.
\end{proposition}

\begin{tcolorbox}[colback=blue!2,colframe=blue!50!black,arc=2mm,
  title={\textbf{The operator, specified as far as honestly possible}},
  fonttitle=\bfseries,breakable]
The missing self-adjoint operator is now pinned:
\begin{itemize}
\item \textbf{class:} $\C$ / unitary / $\beta=2$ (\emph{not} the substrate's
      $\mathbb{H}$/$\beta=4$);
\item \textbf{finite-place part:} the verified commutative Hecke (scaling)
      algebra (\S\ref{sec:locate});
\item \textbf{what is missing (\TAGopen):} the \emph{archimedean} dilation
      glued in, with a \emph{self-adjoint} realisation making the zeros a
      discrete point spectrum --- i.e.\ Connes' Weil positivity for
      \emph{all} test functions. That last step is RH.
\end{itemize}
\end{tcolorbox}

\section{Status ledger}

\begin{center}
\begin{tabular}{p{0.66\textwidth}l}
\toprule
Claim & Status\\
\midrule
$120$ icosian units $=2I$; norms, closure & \TAGver\\
Each local factor's zeros on $\Re(s)=1/2$ (Satake circle) & \TAGver\\
Substrate computes the genuine norm-$31$ newform ($11$ ideals, out-of-sample) & \TAGver\\
Multiplicativity / Ramanujan / Sato--Tate on the eigenform & \TAGver/\TAGthm\\
$L(\psi,s)$ is a genuine cuspidal automorphic L-function & \TAGthm\\
Connes/Weil functional $W(h)>0$ on validated tests & \TAGev\\
Zeros are GUE (quantum chaos); diffract into primes (quasicrystal) & \TAGver\\
$\mathbb{H}\to\C$ class flip $=$ complexification ($2I\subset\mathrm{SU}(2)$) & \TAGver\\
Symmetry type: bulk Unitary/GUE, family Orthogonal $O^+$ & \TAGver/\TAGthm\\
Adelic scaling: finite-place Hecke algebra is commutative, $\C$-class & \TAGver\\
``Admissibility $+$ explicit formula $\Rightarrow$ RH'' & withdrawn (false)\\
RH for $L(\psi,s)$ (= GRH); the archimedean self-adjoint operator & \TAGopen\\
Existence of the Hilbert--Pólya operator $\hat H$ & \TAGconj/\TAGopen\\
\bottomrule
\end{tabular}
\end{center}

\section{Scope discipline}

The mathematics above uses only fields, lattices, quaternion orders, Hecke
operators, gamma factors, and conductors. \textbf{No observer, no
consciousness, and no ``self'' enters any definition, and none is needed.}
Any reading of primes or zeros as selves/agents is a personal-philosophical
frame that is \emph{outside} this mathematics; the work neither uses nor
supports it, consistent with the programme's own discipline (``$X$ exists
$\Leftrightarrow \mathcal{C}(X)=X$''; no ``primes are alive''; no ``universe
is conscious''). The randomness of the primes is the friction between
addition and multiplication --- field-independent, sourced by no value or
agent; RH is the timeless, observer-free statement that the primes are
maximally balanced.

\section*{Conclusion}
The substrate reaches the RH frontier exactly this far: it \emph{verifiably
computes} a genuine Hilbert modular newform from its own geometry, realises
the \emph{local} critical line in full, supplies real Satake data to the
genuine Connes positivity object (positive on every validated test), and ---
through the quantum-chaos and division-algebra structure --- pins the missing
operator to the $\C$/unitary/$\beta=2$ class, with its finite-place part
(the commutative Hecke scaling algebra) built and verified. It does
\emph{not} prove RH: the single remaining step is the \emph{archimedean
self-adjoint realisation} of the zeros as a discrete spectrum --- equivalently
Weil positivity for all test functions, equivalently the Hilbert--Pólya
operator on the $\C$/adelic side --- which is RH and is open for everyone.
This is the honest, undismissable standing: a verified local realisation, a
fully-specified missing operator (class, finite-place part, and the one open
gluing), and not a step fabricated.

\begin{thebibliography}{9}
\bibitem{Dembele} Dembl\'e, L. \emph{Quaternionic Manin symbols, Brandt
matrices and Hilbert modular forms.} Math.\ Comp.\ 76 (2007), 1039--1057.
\bibitem{Voight} Voight, J. \emph{Quaternion Algebras.} GTM 288, Springer, 2021.
\bibitem{Connes} Connes, A. \emph{Trace formula in noncommutative geometry
and the zeros of the Riemann zeta function.} Selecta Math.\ 5 (1999), 29--106.
\bibitem{BerryKeating} Berry, M.V., Keating, J.P. \emph{The Riemann zeros and
eigenvalue asymptotics.} SIAM Review 41 (1999), 236--266.
\bibitem{IK} Iwaniec, H., Kowalski, E. \emph{Analytic Number Theory.} AMS
Colloq.\ Publ.\ 53, 2004.
\bibitem{Montgomery} Montgomery, H.L. \emph{The pair correlation of zeros of
the zeta function.} Proc.\ Sympos.\ Pure Math.\ 24 (1973), 181--193.
\bibitem{RudnickSarnak} Rudnick, Z., Sarnak, P. \emph{Zeros of principal
L-functions and random matrix theory.} Duke Math.\ J.\ 81 (1996), 269--322.
\bibitem{Atas} Atas, Y.Y., Bogomolny, E., Giraud, O., Roux, G.
\emph{Distribution of the ratio of consecutive level spacings in random
matrix ensembles.} Phys.\ Rev.\ Lett.\ 110 (2013), 084101.
\bibitem{Dyson} Dyson, F.J. \emph{The threefold way.} J.\ Math.\ Phys.\ 3
(1962), 1199--1215.
\bibitem{KatzSarnak} Katz, N.M., Sarnak, P. \emph{Random Matrices, Frobenius
Eigenvalues, and Monodromy.} AMS Colloq.\ Publ.\ 45, 1999.
\bibitem{ILS} Iwaniec, H., Luo, W., Sarnak, P. \emph{Low lying zeros of
families of L-functions.} Publ.\ IH\'ES 91 (2000), 55--131.
\end{thebibliography}

\end{document}
